\section{Cross-Receiver Failure Diagnosis}
\label{sec:diagnosis}

\subsection{Problem formulation}
We consider closed-set LoRa device identification under cross-receiver transfer.
A source receiver $r_s$ provides labeled training captures for $C$ devices.
A target receiver $r_t$ provides unlabeled calibration traffic and a held-out query set for evaluation.
The goal is to classify each target-receiver device file correctly after minimal target-side calibration.

We use a frozen hybrid RF feature extractor trained on the source receiver~\cite{paper1_placeholder}.
All diagnosis in this section is performed on extracted embeddings and OOB spectra; we do not retrain the backbone.

\subsection{Receiver-induced feature entanglement}
Cross-receiver failure is not explained by complete embedding collapse alone.
File-level cosine analysis shows that same-device cross-receiver pairs can remain relatively close for the hybrid fused embedding, while device separability and decision boundaries degrade.
We therefore analyze receiver entanglement at three levels: physical OOB spectra, linear probes on embedding paths, and prediction collapse.

\subsection{OOB path receiver probe}
We train linear receiver probes on file-level mean embeddings mixed across RX1 and RX2.
For the hybrid model, the OOB-only path yields the highest receiver-probe accuracy ($\approx 0.73$), compared with main ($\approx 0.62$) and fused ($\approx 0.50$).
This indicates that OOB evidence, although useful under fixed-receiver operation, becomes more receiver-entangled under cross-receiver shift.

\subsection{OOB spectral receiver bias}
Physical-layer OOB analysis shows a systematic RX2/RX1 OOB energy ratio of $1.44$, supporting receiver front-end spectral bias rather than transmitter-only variation.

\subsection{Embedding distance ratio}
File-level cosine distance ratios compare same-device cross-receiver similarity against different-device same-receiver similarity.
CNN-IQ exhibits ratio $> 1$, indicating receiver shift dominates device separation.
The hybrid fused embedding shows ratio $< 1$, meaning same-device cross-receiver structure is partially preserved, but device probe accuracy on mixed receivers remains low ($\approx 0.29$ for 24 classes).

\subsection{Confusion and prediction collapse}
Under RX1$\rightarrow$RX2, CNN-IQ collapses to one or two dominant predicted classes (top-1 mass up to $95.8\%$).
The hybrid model performs better but remains near chance for 24-way file-level classification.
This motivates labeled target calibration rather than unsupervised pseudo-label adaptation alone.

\subsection{Diagnosis conclusion}
Cross-receiver failure arises from receiver-induced OOB entanglement and non-transferable source decision geometry.
The frozen hybrid backbone retains partial same-device structure, making lightweight target-receiver prototype calibration a plausible repair mechanism.

\begin{figure}[t]
  \centering
  \includegraphics[width=\linewidth]{fig1_diagnosis_summary.pdf}
  \caption{Diagnosis summary: (a) receiver probe by path, (b) OOB spectral bias, (c) distance ratio, (d) collapse.}
  \label{fig:diagnosis}
\end{figure}
